\subsection{Looped Transformers}

In this section we introduce Looped Transformers and define the notation that we will use throughout the paper. We represent the input sequence to our transformer of length $T$ and dimension $D$ as $\mX \in \mathbb{R}^{T \times D}$. Following the notation of \citet{yudin2025pay} we define $\attn$, the dot product self-attention mechanism, as a map $f: \mathbb{R}^{T \times D} \to \mathbb{R}^{T \times D}$:
\begin{align}
    \attn(\mX) &= \softmax{\left ( \frac{\mX \mW_Q \mW_K^{\top} \mX^{\top}}{\sqrt{d}} \right )} \mX\mW_V, \\ 
    &= \softmax{\left ( A(\mX) \right )} \mX\mW_V, \label{eq:attn_in_a}
\end{align}
where $\mW_Q, \mW_K, \mW_V \in \mathbb{R}^{D \times d}$ are projection matrices and $A$ is defined for convenience. 

A transformer block typically comprises an attention mechanism and a position-wise MLP as follows:
\begin{align}
    \hat{\mX} &= n_2 \left ( \mX + \attn(n_1(\mX) \right ), \label{eq:block_norm_1} \\
    {\mX}' &= n_4 \left ( \hat{\mX} + \text{MLP}(n_3(\hat{\mX})) \right ), \label{eq:block_norm_2}
\end{align}
where $n_1, n_2, n_3, n_4$ are each optional norms -- here we are borrowing from the notation of \citet{geiping2025scaling}. We denote the action of a Transformer block $\block: \mathbb{R}^{T \times D} \to \mathbb{R}^{T \times D}$ as $\mX' = \block(\mX)$, and refer to the intermediate hidden-state matrices $\mX$ between blocks as the \textbf{residual stream}.
% We will denote the action of this Transformer block $\block: \mathbb{R}^{T \times D} \to \mathbb{R}^{T \times D}$ as
% \begin{equation}
%     {\mX}' = \block(\mX),
% \end{equation}

% and the matrices $\mX$ of hidden states between Transformer blocks we will frequently refer to as the \textbf{residual stream}.

Looped Transformers are Transformers that utilize ``recurrence in depth'' -- that is, they reapply layers to repeatedly act on the latent states. Recent research has identified that an effective way \citep{bae2025mixture} to achieve this is via \emph{cyclic recurrence}: a fixed sequence of layers is repeated in a ``cyclic'' pattern. This is the approach that we focus on in this work, and we introduce it in more detail below. For convenience, we define a $k$-stacked block as a composition of Transformer blocks,
$\stack_k(\mX) = \block_k(\block_{k-1}(\dots \block_1(\mX) \dots))$.


% For convenience, we define a $k$-stacked block  comprising a sequence of separate Transformer blocks:
% \begin{equation}
%     \stack_k(\mX) = \block_k(\block_{k-1}( \dots \block_1(\mX) \dots ))
% \end{equation}

In the case of \citet{geiping2025scaling,mcleish2025teaching} this stacked block may also take an additional input $\mZ \in \mathbb{R}^{T \times D}$ which is typically initialized from a Normal distribution, as well as the original input to the recurrent section: this is known as \emph{input injection} \citep{bai2019deep,anil2022path}, and the two inputs are projected into common feature space $\mathbb{R}^D$ before the block is applied. In the case of input injection, a $k$-stacked block therefore becomes
\begin{equation}
    \stack_k(\mX, \mZ) = \block_k ( \block_{k-1}( \dots \block_1([\mX, \mZ] \mW_I) \dots )),
\end{equation}
where $[\cdot, \cdot]$ denotes concatenation in the channel dimension and $\mW_I \in \mathbb{R}^{2D \times D}$ is a learned projection matrix.

This allows us to define a $(k \otimes l)$-Recurrent block as a $k$-stacked block repeated $l$ times:
\begin{equation}
    R_{l,k}(\mX) = \overbrace{\stack_k ( \stack_{k}( \dots \stack_k(}^{\times l} \mX) \dots )),
\end{equation}

which with input-injection becomes
\begin{equation}
    R_{l,k}(\mX, \mZ) = \overbrace{\stack_k ( \mX, \stack_{k}( \dots \mX, \stack_k( \mX, \mZ)) \dots ))}^{\times l}.
\end{equation}

Note that the input $\mX$ is only ``injected'' at the start of each stack of blocks -- once per recurrence. Additionally, a complete looped Transformer may have multiple feedforward layers before the Recurrent block, and multiple feedforward layers after the Recurrent block: following the convention of \citet{geiping2025scaling}, we refer to these as \emph{prelude} and \emph{coda} layers respectively, and these are simply separate stacked blocks with non-tied layer weights. Where prelude and coda are used, we will frequently refer to this as a \emph{sandwich} block structure.

We combine and adapt the nomenclature of \citet{geiping2025scaling,saunshi2025reasoning} and refer to a looped Transformer with $p$ prelude layers, $k$ recurrent layers and $c$ coda layers with the tuple $(p, k, c)$. Where input injection is used, we add an $I$ subscript $(p, k, c)_I$, and when referring to a recurrent layer looped a specific number of times $l$, we denote this as $(p, k \otimes l, c)$.

In summary, a $(p, k \otimes l, c)$ looped Transformer is defined as
\begin{align*}
    \mX_0 &\gets \stack_p(\mX) \\
    \mX_i &\gets \stack_{k}^{\prime}(\mX_{i-1}) \qquad i \in \{1, \dots, l \} \\
    \mX &\gets \stack_c^{\prime \prime}(\mX_l),
\end{align*}
where $\prime$ and $\prime \prime$ indicates that these are \emph{different} stacks between which weights are not shared. A $(p, k \otimes l, c)_I$ looped Transformer (with input injection) is defined as
\begin{align*}
    \mX &\gets \stack_p(\mX) \\
    \mZ_i &\gets \stack_{k}^{\prime}(\mX, \mZ_{i-1}) \qquad i \in \{1, \dots, l \} \\
    \mX &\gets \stack_c^{\prime \prime}(\mZ_l),
\end{align*}
where $\mZ_0$ is initialized such that each column is sampled from $\mathcal{N}(\boldsymbol{0}, \sigma^2 \mathbb{I}_{D})$.

% We will often refer to the depth of the entire looped model -- summing the depth of stacked layers -- as the \textbf{realized depth}, and the depth of the cyclic recurrent block as the \textbf{recurrent depth}.

We will describe our results via \textbf{grouping}:
\begin{itemize}
    \item \textbf{No grouping}: The value is visualized as it evolves through sequential layers of the model, irrespective of whether these layers are repeated.
    \item \textbf{By recurrence}: Separate lines are visualized for each complete pass through the recurrent block. The $x$-axis is typically percentage scaled to represent relative depth within that block (including prelude/coda), allowing us to overlay and compare successive passes. Always presented with a \textbf{\textcolor{mplsummergreen}{green}-\textcolor{mplsummeryellow}{yellow}} colorbar, with later recurrences colored more yellow.
    \item \textbf{By layer}: Separate lines are visualized for each unique layer, showing how the value evolves across recurrences. Always presented with a \textbf{\textcolor{mplwinterblue}{blue}-\textcolor{mplwintergreen}{green}} colorbar, with later layers colored more green.
\end{itemize}

\subsection{Stages of Inference}\label{subsec:stages_of_inference_prelim}

The behavior of layers in feedforward Transformers appears to change sharply with depth: \citet{lad2024remarkable} originate the term ``stages of inference'' and demonstrate how several different layer mechanisms emerge at different depths. \citet{queipo2025attention} further develop this viewpoint, focusing on behaviors that can be characterized by the \emph{mixing} (or lack thereof) induced by the attention heads. We focus on this latter perspective.
% For a representation matrix $\mathbf{X}$ with singular values $\sigma_1 \geq \sigma_2 \geq \ldots \geq \sigma_r$, we measure compression via the \textbf{matrix-based entropy}:
% \begin{equation}
% H(\mathbf{X}) = -\sum_{j=1}^r p_j \log p_j, \\ \text{where } p_j = \sigma_j^2/\|\mathbf{X}\|_F^2
% \end{equation}
% Low entropy indicates compression into few dominant directions. The \textbf{anisotropy} $p_1 = \sigma_1^2/\|\mathbf{X}\|_F^2$ measures directional bias~\citep{razzhigaev2024shapelearninganisotropyintrinsic} (near 1 = extreme bias, near $1/r$ = isotropy).
Mixing in this context refers to the extent to which the attention mechanism incorporates information from previous tokens at each layer. Throughout the main text of this paper we quantify our study of mixing behavior through the \emph{ColSum Concentration} metric introduced in \citet{queipo2025attention}; we introduce and discuss additional metrics in \Cref{appen:additional_soi_results}.

We first define the column sum $c_j = \sum_i A_{ij}$ to capture how much attention mass is received by token $j$. We normalize this as $\hat{c}_j = c_j / T$ to obtain a probability distribution, noting that $\sum_{i,j} A_{ij} = T$ since $A$ is row-stochastic. From this distribution, we define the \textbf{ColSum Concentration} via its normalized entropy as
$C = 1 - H_{\text{col}} \in [0,1] = 1 + \frac{1}{\log T} \sum_j c_j \log c_j$.
 

% From this distribution we define the \textbf{ColSum Concentration} from its normalized entropy:
% \begin{align}
%     C &= 1 - H_{\text{col}} \in [0,1] \\
%     &= 1 + \frac{1}{\log{T}} \sum_j c_j \log{c_j}
% \end{align}
Large values of $C$ indicate a high \emph{concentration} of attention mass: few columns receive most of the mass. We note therefore that this metric captures the well-studied \emph{attention sink} \citep{xiao2023efficient,barbero2025llms} behavior, but generalizes to capture concentration over any token position. This property is useful for our investigation as not all models studied herein exhibit sinks on the first token; in particular OLMo-2 frequently concentrates attention mass on punctuation, echoing a result in \citet{sandoval2025using}.